\section{Conclusion}
\label{sec:conclusion}
This revision re-examined E-DDQN, a dueling Double DQN with an asymmetric (5:1) reward, class-balanced replay and soft target updates, for safety-prioritised cybersickness detection motivated by an adversarial threat model, and compared it with a standard DQN and with cost-sensitive supervised models. Rebuilding the pipeline from the original UFFCSData exports revealed that the previously reported performance was produced by oversampling before a row-level split; with the same data and a grouped split, performance falls to near chance for the original columns. We corrected the E-DDQN implementation (valid within-session transitions with terminal flags, removal of an invalid teacher-forcing scheme, target-network normalisation statistics), proved that on logged data the problem is a contextual bandit whose asymmetric-reward optimum equals the 5:1 cost-sensitive Bayes threshold, and evaluated all models under one frozen, nested, participant-conservative protocol. Under that protocol E-DDQN reached cost \EDDQNCostCostMean{} and ROC-AUC \EDDQNCostRocAucMean, did not outperform the standard DQN (cost \DQNCostCostMean, ROC-AUC \DQNCostRocAucMean), and was significantly worse than a cost-weighted random forest (cost \RFCostCostMean, ROC-AUC \RFCostRocAucMean); only the random forest improved materially on always alarming (cost \AlwaysCost), and with \nGroupsEval{} independent groups even that improvement is not statistically established. Bounded telemetry perturbations changed threshold-free performance by at most \AttAucMaxAbsChange{} ROC-AUC, while adaptive worst-case manipulation reduced recall on event seconds. The safety objective of near-complete detection with controlled false alarms is not achieved on these data. Future work requires datasets with physiological or eye-tracking signals, verified participant identifiers and, for the adversarial claim, manipulations applied to participants with measured effects.
